%% ===========================================================================
\FloatBarrier
\section{Finite element formulation}
\label{sec:fem}
%% ===========================================================================

\FloatBarrier
\subsection{Weak form}
\label{sec:weak}

Testing \eqref{eq:gov1} and \eqref{eq:gov2} against a virtual displacement
$\delta\bm{u}$ and a virtual potential $\delta\phi$ and integrating by parts
twice, the strain-gradient term produces both the double traction
\eqref{eq:doubletraction} and an integral over the element interiors,
\begin{multline}
  \int_\Omega \delta\varepsilon_{ij}\,C_{ijkl}\,\varepsilon_{kl}\,\mathrm{d}A
  + \tfrac12\int_\Omega \delta\varepsilon_{ij,m}\,M_{mn}\,C_{ijkl}\,
     \varepsilon_{kl,n}\,\mathrm{d}A \\
  - \int_\Omega \delta\varepsilon_{ij}\,e_{kij}\,E_k\,\mathrm{d}A
  - \int_\Omega \delta\varepsilon_{ij,m}\,\mu_{kijm}\,E_k\,\mathrm{d}A \\
  - \int_\Omega \delta E_k\left(e_{kij}\varepsilon_{ij}
     + \kappa_{kl}E_l + \mu_{kijl}\eta_{ijl}\right)\mathrm{d}A
  = \delta W_{\rm ext}
  + \int_{\Gamma}\left(\delta u_i\,t_i + \delta u_{i,j}n_j\,m_i\right)
    \mathrm{d}s ,
  \label{eq:weak}
\end{multline}
which is the form that the element matrices implement. The boundary integral
carries the two mechanical conditions of Sec.~\ref{sec:bc_mech}; the electrical
term carries the electrode conditions of Sec.~\ref{sec:bc_elec} and is
condensed before assembly, so that the mechanical problem never sees the
potential degrees of freedom.

\FloatBarrier
\subsection{The BFS element and its degrees of freedom}
\label{sec:bfs}

The domain is covered by Bogner--Fox--Schmit rectangles. Each element has four
nodes and, per node, the eight mechanical degrees of freedom \eqref{eq:dofs}
carried by a bicubic Hermite tensor-product interpolation, that is 32
mechanical degrees of freedom per element, together with the four nodal
potentials of a bilinear interpolation. The strain and the strain gradient on
the element are
\begin{equation}
  \varepsilon = B(\xi,\eta)\,\bm{u}^e,\qquad
  \eta_k = \partial_k\varepsilon = B_{,k}(\xi,\eta)\,\bm{u}^e,
  \label{eq:Bmat}
\end{equation}
both integrated by $2\times2$ Gauss quadrature in the element frame; the
gradient matrix $B_{,k}$ contains the third derivatives of the shape functions
and is the reason the element has to be $C^1$. The element contributions are
the elastic $K_e=\int B^{\mathsf T}CB$, the gradient
$K_{g,e}=\tfrac12\int M_{mn}B_{,m}^{\mathsf T}CB_{,n}$, the piezoelectric coupling
$B_e=\int B^{\mathsf T}e^{\mathsf T}\nabla N_\phi$, the dielectric
$C_{\phi,e}=-\int\kappa\,\nabla N_\phi\cdot\nabla N_\phi$ and a consistent mass
matrix from the same interpolation. The flexoelectric contribution is
assembled into $B_e$ as the additional potential--displacement coupling that
the $\mu_{kijl}\eta_{ijl}$ term of \eqref{eq:energy} produces; because it
involves $B_{,k}$, it is the term for which the $C^1$ character of the element
is indispensable.

\FloatBarrier
\subsection{Condensation of the potential}
\label{sec:condensation}

The assembled system is indefinite in the saddle-point form
\begin{equation}
  \begin{bmatrix} K & B \\ B^{\mathsf T} & -C_\phi \end{bmatrix}
  \begin{bmatrix} \bm{u} \\ \bm{\phi}\end{bmatrix}
  = \begin{bmatrix} \bm{f} \\ \bm{q}\end{bmatrix},
  \label{eq:saddle}
\end{equation}
with $C_\phi$ positive by construction. Eliminating the potential degrees of
freedom of the electrode set reproduces \eqref{eq:rankone} exactly, and the
returned pair $(d,\beta)$ is what the modal and the static solvers consume.
Two conventions are enforced in the implementation because they are easy to get
wrong: the charge is read as $q=d^{\mathsf T}\bm{u}/\beta$ rather than by
pointwise integration of the electric displacement, and eigenvectors are
phase-normalized on a nodal displacement degree of freedom -- normalizing on
$\max|\phi|$ can select a derivative degree of freedom and return a
meaningless voltage.

\FloatBarrier
\subsection{Modal L1 time stepping}
\label{sec:stepping}

With the L1 weights of \eqref{eq:l1} the time-discrete problem is
\begin{equation}
  \left(M + \Delta t^{\alpha}\Gamma(2-\alpha)\,K_{\rm mem}\right)\bm{a}^n
  = \bm{f}^n - K\bm{u}^n - K_{\rm mem}\sum_{k=1}^{n-1} b_{n-k}
    \left(\bm{u}^{k}-\bm{u}^{k-1}\right),
  \label{eq:step}
\end{equation}
solved by a Cholesky factorization of the constant left-hand matrix, with
$K_{\rm mem}\in\{K_{\rm eff},K_{uu}\}$ the model choice quantified in
Sec.~\ref{sec:memorystudy}. In the modal variant used for the long ring-downs,
the same convolution is carried on the mass-normalized coordinates of the
retained modes, so the cost per step is proportional to the number of retained
modes rather than to the number of degrees of freedom; the truncation error of
that projection is measured against the full-order solution rather than
assumed small (Sec.~\ref{sec:verification}).

\FloatBarrier
\subsection{Implementation remarks}
\label{sec:implementation}

Two properties of the implementation are worth recording, because they
govern both cost and accuracy.

\begin{enumerate}
\item \emph{Scaling.} The gradient block is $O(M/h^4)$ while the elastic block
  is $O(C/h)$, so the assembled pencil becomes badly scaled as the thickness
  falls. The implementation scales the two blocks separately before the
  eigensolve and returns the eigenpairs unscaled; the thin-plate limit of
  Sec.~\ref{sec:thickness} is at a condition number of about $10^{12}$, and the
  sensitivity of the resulting small differences is measured rather than
  assumed (\ref{app:verification}).
\item \emph{Cost.} A mesh of $48\times4$ BFS elements has $49\times5=245$
  nodes and $8\times245=1960$ mechanical degrees of freedom and is the
  workhorse of the results; the mesh-refinement study of
  Sec.~\ref{sec:verification} shows that on that mesh the frequency of the
  extensional mode is within $0.05\,\%$ and its quasi-static coupling factor
  within $0.03\,\%$ of a $96\times4$ refinement, and the thinnest plates are
  solved on it because the conditioning, not the number of elements, is the
  binding constraint.
\end{enumerate}

\FloatBarrier
